% ======================================================================
% Chapter 2 -- Part B: writing the design down (step 4); patterns and
% anti-patterns
% ======================================================================

\section{Step 4: How Do We Write a Design Down?}
\label{sec:design-table}
\label{sec:notation}

\begin{inpractice}[In practice: three engineers, one design, three stories]
After Harbor funded Team B's quote design, three engineers described it on the company wiki. One wrote that ``the agent looks up pricing and checks policy.'' Another wrote that ``a rules engine applies the discount cap.'' The third wrote that ``the model decides whether a manager needs to approve.'' All three were describing the same system, and the third description was wrong, because in Team B's design code decides when a manager is needed. Nobody could tell from the prose, because words such as \textbf{agent}, \textbf{tool} and \emph{checks} meant different things to each writer. The first design review took two hours, most of which were spent agreeing on what the design was before anyone could ask whether it was good.
\end{inpractice}

The remedy is not more prose but a fixed format that forces the same few facts onto the page for every part of a design. Step 4 of the method asks you to draft two or three designs and to write each of them down in this format.

\begin{definition}
A \keyterm{design table}\index{design table|textbf} lists the steps of a solution in the order in which the work flows, and states for each step who or what does it, what it uses, and why it is there.
\end{definition}

Each of the four columns has a simple rule.
\begin{description}[leftmargin=1em,itemsep=2pt]
  \item[Step.] Write a short phrase that says what happens, such as ``check the warranty.'' When the work can go more than one way, write every branch into the step, including the ``otherwise'' case, for example ``if any field is missing, ask the representative; otherwise continue.''
  \item[Done by.] Write one of four kinds of worker, \tagPerson, \tagAgent, \tagCode\ or \tagGuardrail, and add the one fact a reviewer needs about it. For a \textbf{person}, give the role and how quickly they must respond. For an \textbf{agent}, say how much it may decide and when it stops, for example ``one model call; stops when the form is filled.'' For \textbf{code}, say where the rule comes from, such as ``configuration version 3.'' For a \textbf{guardrail\index{guardrail}}, an automatic check that blocks or repairs anything that breaks a policy, say what it checks.
  \item[Uses.] Write the MCP servers, knowledge stores and inputs that the step relies on. For an MCP server, give its owner and list only the tools this solution may use; that list is the same allow-list you will give ADK\index{Agent Development Kit (ADK)}'s \texttt{tool\_filter} in code (Listing~1.4), and it is where most security decisions live. For a knowledge store, give its owner and how often it is refreshed. A company system such as the ERP is always reached through a server, never directly, so it appears here only as the server's name.
  \item[Why it is there.] Write the business measure the step improves or the risk it guards against. This column enforces a simple rule: \emph{every step must earn its place}. A step with nothing in this column either is missing its reason or should not be there.
\end{description}

Table~\ref{tab:quote-design} writes Harbor's quote design from Chapter~\ref{ch:foundations} in this format. The last column refers to three business measures and three risks. The measures are \textbf{turnaround}, the median time to quote, which should fall below 4 hours from 26; \textbf{margin}, the average realized margin, which should stay at least at today's level; and \textbf{analyst time}, the pricing analyst's minutes per quote, which should fall below 8 from about 35. The risks are a \emph{below-floor} quote issued under the minimum margin, the \emph{wrong customer} or contract applied to a quote, and a \textbf{mismatch} between a number in the cover note and the quote itself.

\begin{table}[htbp]
\centering
\small
\caption{Harbor's quote design as a design table. It is the same design as the sketch in Chapter~\ref{ch:foundations}, but the kind, owner, rule source and reason for every step are now on the page.}
\label{tab:quote-design}
\begin{tabularx}{\textwidth}{@{}c>{\raggedright\arraybackslash}X>{\raggedright\arraybackslash}p{3.0cm}>{\raggedright\arraybackslash}p{3.3cm}>{\raggedright\arraybackslash}p{2.0cm}@{}}
\toprule
\rowcolor{coHead} \# & Step & Done by & Uses & Why it is there \\
\midrule
1 & Screen the representative's request, which arrives by e-mail or chat and is untrusted text & \tagGuardrail\ Blocks instructions aimed at the system & The request & Wrong customer \\
2 & Read the request and fill in the quote form & \tagAgent\ One model call; stops when the form is filled & The request & Turnaround; analyst time \\
3 & Validate the fields and fetch prices. If a field is missing, ask the representative; otherwise continue & \tagCode & ERP server (owner: Pricing), tools \texttt{get\_item\_pricing} and \texttt{get\_contract\_terms} & Wrong customer \\
4 & Apply the discount policy. If the quote is an exception, send it to step 5; otherwise go to step 6 & \tagCode\ Discount cap 15\%, margin floor 22\%, from configuration version 3 & Policy and contract index (owner: Finance; refreshed daily) & Margin; below floor \\
5 & Approve or reject the exception. A rejected quote returns to the representative & \tagPerson\ Sales manager, within 4 business hours & The quote and the reason it is an exception & Below floor \\
6 & Draft the cover note & \tagAgent\ One model call; stops when the note is written & The priced quote from step 4 & Analyst time \\
7 & Check every number in the cover note against the policy engine's output & \tagGuardrail & The note and the priced quote & Mismatch \\
8 & Save the quote as a draft in the CRM & \tagCode & CRM server (owner: Sales Ops), tool \texttt{create\_quote\_draft} & Turnaround \\
9 & Review the draft and send it to the customer & \tagPerson\ The sales representative, who now owns the quote & The CRM draft & Turnaround \\
\bottomrule
\end{tabularx}
\end{table}

Three things are now explicit that the informal sketch left vague. The model never decides whether a manager is needed, because step 4 is done by code. Both agents have bounded autonomy, since each makes a single call and has a defined stopping point. And the one step that changes a company system, saving the draft in step 8, comes after a guardrail that checks every number in the cover note, and is followed by a person who reviews the quote before it reaches the customer.

\subsection{Eight checks before a review}
\label{sec:checks}

A design table is \emph{complete} when it passes the eight checks below. They are mechanical on purpose, so that a reviewer can apply them before discussing whether the design is any good, and each one catches a mistake that recurs in student designs and in production systems alike. Each check is a question you can answer by reading down one column of the table.

\begin{tryit}[Checklist: eight checks before a review]
\begin{description}[leftmargin=1em,itemsep=2pt]
  \item[$\square$ Bounded agents.] Does every agent state how much it may decide and when it stops? An agent with no stopping point has no limit on its cost.
  \item[$\square$ Protected writes.] Does every step that changes a company system, sends a message outside the company or moves money come after a person or a guardrail, or carry a note explaining why neither is needed?
  \item[$\square$ Servers in front of systems.] Does every step reach company systems only through an MCP server with a listed set of tools?
  \item[$\square$ Rules in code.] Are thresholds, limits and eligibility rules done by code rather than written into an agent's instructions?
  \item[$\square$ Owned knowledge.] Does every knowledge store name an owner and say how often it is refreshed?
  \item[$\square$ Screened input.] Is every piece of untrusted content, such as e-mail from outside the company or a document uploaded by a customer, screened by a guardrail before an agent reads it, or is there a note explaining why it need not be?
  \item[$\square$ Complete branches.] Does every step that can go more than one way say what happens in every case, including ``otherwise''?
  \item[$\square$ Justified steps.] Does every step have an entry in the ``why it is there'' column?
\end{description}
\end{tryit}

Complete is not the same as good. A complete design can still be too costly, too slow, or aimed at the wrong problem. The checks guarantee only that the design is explicit enough to be judged, and the judging is the subject of Section~\ref{sec:criteria}.

\paragraph{Component cards.} A table cannot hold every detail, so each agent, MCP server and guardrail also gets a short written specification called a \keyterm{component card}\index{component card|textbf}. The card states the component's purpose, its inputs and outputs, its autonomy and stopping point, its tools, a quality target, a cost and time budget, and what happens when it fails. The quality target and the budget are what make a design testable. The first case, Chapter~\ref{ch:case-policy}, shows a complete card, and Chapter~\ref{ch:case-fraud} shows how to build the test set that a quality target refers to.

% ----------------------------------------------------------------------
\section{Patterns and Traps: Do We Have to Start from Scratch?}
\label{sec:patterns}

When you draft designs in step 4, you rarely need to start from nothing. A \keyterm{design pattern}\index{design pattern|textbf} is a named, reusable arrangement of steps that solves a recurring problem, with known costs and known ways to fail. Naming patterns lets a team say ``put an approval gate in front of the credit'' instead of describing the arrangement from scratch each time. This chapter presents only the two patterns that appear in almost every case: \textbf{retrieve-then-answer}, which brings company knowledge in, and the \textbf{approval gate}\index{approval gate|textbf}, which keeps a person in charge of consequential actions. Other patterns are introduced in the chapter that first needs them.

% ---------------- retrieve-then-answer
\begin{patterncard}{Retrieve-then-answer}\index{retrieve-then-answer|textbf}
\textbf{Intent.} Answer from the company's own documents rather than the model's memory, and cite the source. This is the RAG\index{retrieval-augmented generation (RAG)} building block of Chapter~\ref{ch:foundations} used as a pattern.\par\smallskip
\textbf{How it works.} Code rewrites the user's question and searches an index of company documents. An agent then answers using only the passages found, and cites each one. Finally, a guardrail checks that every claim in the answer has a citation, and an answer without one is not sent.\par\smallskip
\textbf{Use when} answers depend on private or changing text and users must be able to check the source. \textbf{Avoid when} the answer is a single database field, which should be looked up instead. \textbf{Cost} is one search and one model call per question, which is low and predictable. \textbf{Typical failures} are a right answer from an outdated document, a plausible answer with no supporting passage, and leaking documents the user may not see. \textbf{In this book:} Chapter~\ref{ch:case-policy}, and the knowledge parts of Chapters~\ref{ch:case-incident} and~\ref{ch:case-loan}.
\end{patterncard}

% ---------------- approval gate
\begin{patterncard}{Approval gate}
\textbf{Intent.} Pause before a consequential action until an authorized person approves, edits or rejects it.\par\smallskip
\textbf{How it works.} An agent or a code step proposes an action. Code then decides, by a written rule, whether the action needs approval. If it does, the action waits until an authorized person approves, edits or rejects it; if it does not, it goes ahead at once. Only after that does code call the server tool that makes the change.\par\smallskip
\textbf{Use when} an action is costly or hard to undo, or policy assigns the decision to a role. \textbf{Avoid when} the action is cheap and reversible; there a gate only adds delay, and a review of a sample after the fact is better. \textbf{Cost} is mostly the approver's time and the wait, both of which belong in the business case. \textbf{Typical failures} are rubber-stamping\index{rubber-stamping} when approvals are too frequent, requests piling up when the approver is away, and approvers shown too little context to judge. \textbf{In this book:} Chapter~\ref{ch:case-invoice}, and also Chapters~\ref{ch:case-fraud}, \ref{ch:case-kyc} and~\ref{ch:case-loan}.
\end{patterncard}

Five more patterns appear later, each in the chapter that first needs it. The \textbf{router}, which classifies each input once and sends it down a path built for that kind of input, appears in Chapter~\ref{ch:case-support}. The \textbf{planner--executor}, in which one agent splits an open-ended task among worker agents, and the \textbf{event-triggered agent}, which starts work when something happens in a system rather than when a person asks, appear in Chapter~\ref{ch:case-incident}. \textbf{Reflection}, in which a second step critiques an output and sends it back for revision, appears in Chapter~\ref{ch:case-fraud}, and the \textbf{multi-agent handoff} appears in Chapter~\ref{ch:case-airline}.

\begin{thinkabout}
Before reading the three traps below, look back at Team A's design in Chapter~\ref{ch:foundations}. How many things in it would you call a trap? Which one worried the finance director most?
\end{thinkabout}

An \keyterm{anti-pattern}\index{anti-pattern|textbf} is an arrangement that looks reasonable and fails in a predictable way. Three are common enough to name now.

\begin{antipattern}{The do-everything agent}
One agent with many tools and a long prompt does everything, as Team A's design did. It demos well and fails in production. Its behavior is hard to test because every input can take every path, its permissions are the union of everything it might ever need, and a change to one part of the prompt can change behavior somewhere else. \emph{Remedy:} break the process down with the four questions of step 3, and use fixed workflows\index{workflow} and narrow agents.
\end{antipattern}

\begin{antipattern}{Policy in the prompt}
Thresholds, eligibility rules and prohibitions are written into the model's instructions instead of into code. \emph{Remedy:} move the rule to a code step or a guardrail, and keep one sentence in the prompt that explains the rule so the model behaves consistently with it.
\end{antipattern}

\begin{antipattern}{The silent write}
A write to a company system, a message to a customer or a payment happens with no checkpoint, no guardrail and no log that a person will read. \emph{Remedy:} an approval gate, a guardrail, or at least a reversible write plus a sampled review, chosen according to what a wrong write would cost.
\end{antipattern}
